\section{Varying original-line psh nets on the actual curve-weak numerical space}
\label{sec:native-numerical-curve-weak-psh}

\begin{theorem}[Actual numerical FS functions never take positive infinity]
\label{thm:native-numerical-fs-ne-top}
\lean{EnvelopeProperty.NativeNumericalFubiniStudy.FSData.eval_ne_top}
\leanok
\uses{def:native-numerical-whole-section-fs-data}
Every whole original numerical FS datum has values in
$\mathbf R\cup\{-\infty\}$ at every same original native point.
\end{theorem}
\begin{proof}
Each true section norm is a finite nonnegative real number. Its extended
logarithm is therefore never positive infinity; adding a finite real
weight retains this property. A finite supremum of such terms is still
strictly below positive infinity. Multiplication by the strictly positive
finite reciprocal denominator also retains it, including at negative
infinity.
\end{proof}

\begin{definition}[Curve-weak numerical psh with genuinely varying original lines {\cite[Definition~4.1]{BJ22}, \cite[\S2.2]{FHX}}]
\label{def:native-numerical-curve-weak-psh}
\lean{EnvelopeProperty.NativeNumericalFubiniStudy.IsCurveWeakPSH}
\leanok
\uses{def:native-numerical-whole-section-fs-data,thm:native-numerical-fs-nonvanishing,def:native-numerical-curve-weak-topology}
Let the actual original scheme be integral and projective over its original
field. A function is curve-weak numerical psh of class $\theta$ if it is
upper semicontinuous, not identically $-\infty$, and is the pointwise
infimum of a decreasing net of evaluations of actual nonzero whole
original rational-line FS data, on a nonempty directed preorder. The
actual classes $m_j^{-1}[M_j]$ must converge to $\theta$ in the actual
curve-weak numerical quotient. The actual original lines $M_j$, positive
denominators, whole sections and real weights may all vary.

This has BJ's decreasing varying-rational-line form, using the explicit
curve-weak topology. Agreement with the paper's finite-dimensional real
vector-space topology, and the resulting full general numerical psh
comparison, remain mathematical obligations. No fixed-ray restriction,
ample-approximation theorem or supplied model-comparison certificate is
part of this definition.
\end{definition}

\begin{theorem}[The actual numerical psh range and defining convergence]
\label{thm:native-numerical-curve-weak-psh-properties}
\lean{EnvelopeProperty.NativeNumericalFubiniStudy.IsCurveWeakPSH.exists_ne_bot,EnvelopeProperty.NativeNumericalFubiniStudy.IsCurveWeakPSH.upperSemicontinuous,EnvelopeProperty.NativeNumericalFubiniStudy.IsCurveWeakPSH.ne_top,EnvelopeProperty.NativeNumericalFubiniStudy.isCurveWeakPSH_iff_degreewise}
\leanok
\uses{def:native-numerical-curve-weak-psh,thm:native-numerical-fs-ne-top,thm:native-numerical-curve-weak-convergence}
A curve-weak numerical psh function is upper semicontinuous, not identically
$-\infty$, and nowhere positive infinity. Its actual class convergence
condition can equivalently be tested by convergence of every true full
Euler degree on every actual original closed integral curve.
\end{theorem}
\begin{proof}
Upper semicontinuity and nontriviality are explicit defining conditions.
At every original native point, the infimum is bounded above by the
value of any member of the nonempty approximating net. The preceding
true FS bound makes this value strictly less than positive infinity.
For the last assertion, apply the exact degreewise convergence theorem
to the net of the same actual classes of these same original data;
all sections, functions, decreasing order and pointwise infimum remain
unchanged.
\end{proof}

\begin{theorem}[Every actual fixed-ray psh net lies in its genuine numerical psh class]
\label{thm:native-numerical-curve-weak-psh-from-ray}
\lean{EnvelopeProperty.NativeNumericalFubiniStudy.isCurveWeakPSH_of_ray}
\leanok
\uses{def:native-numerical-curve-weak-psh,thm:native-numerical-fs-from-ray,thm:native-numerical-curve-weak-vector-topology,def:native-line-whole-section-psh}
If the existing actual original fixed-$L$-ray predicate holds at real
coefficient $a$, then the same entire original function is curve-weak
numerical psh of its genuine class $a[L]$. No ample hypothesis is required.
\end{theorem}
\begin{proof}
For each member of the actual fixed-ray FS net, choose its original whole
finite datum. Regard its same original power line and sections as a varying
rational-line datum, retaining its denominator and weights. Its evaluation
is literally unchanged and its true class is $a_j[L]$. The given scalar
convergence $a_j\to a$ implies convergence $a_j[L]\to a[L]$ by the proved
joint continuity of real scalar multiplication on the actual curve-weak
quotient. Nonzero families, original decreasing order and exact pointwise
infimum are retained. The original limit already has its required genuine
nontriviality and upper semicontinuity, proving the claim.
\end{proof}

\begin{theorem}[All genuine original curve degrees determine the target class]
\label{thm:native-numerical-curve-weak-psh-intrinsic}
\lean{EnvelopeProperty.NativeNumericalFubiniStudy.isCurveWeakPSH_iff_of_curveDegrees_eq}
\leanok
\uses{def:native-numerical-curve-weak-psh,thm:native-real-numerical-classes}
If two actual original numerical classes have the same true degrees on
all actual original closed integral curves, their curve-weak numerical
psh predicates agree on the same entire original native space.
\end{theorem}
\begin{proof}
The true original curve degrees jointly separate the genuine numerical
quotient, so the two whole target classes are equal. Substitute this exact
equality in the definition. The permitted varying original-line data and
the entire original limit function remain identical. This does not assert
arbitrary numerical invariance of the distinct older fixed-line FS or
fixed-ray psh predicates.
\end{proof}
